% Generated draft: requires mathematical and source review.
\documentclass[11pt]{article}
\usepackage[margin=25mm]{geometry}
\usepackage{fontspec}
\setmainfont{Latin Modern Roman}
\usepackage{amsmath,amssymb,mathtools,unicode-math}
\setmathfont{Latin Modern Math}
\usepackage{enumitem,needspace,xurl,hyperref}
\setlength{\parindent}{0pt}
\setlength{\parskip}{5pt}
\setlength{\emergencystretch}{3em}
\newcommand{\N}{\mathbb{N}}
\newcommand{\Z}{\mathbb{Z}}
\newcommand{\Q}{\mathbb{Q}}
\newcommand{\R}{\mathbb{R}}
\newcommand{\C}{\mathbb{C}}
\newcommand{\F}{\mathbb{F}}
\newcommand{\A}{\mathbb{A}}
\newcommand{\PP}{\mathbb{P}}
\newcommand{\E}{\mathbb{E}}
\newcommand{\eps}{\varepsilon}
\newcommand{\dd}{\,\mathrm{d}}
\begin{document}
\section*{UnsolvedMath ID: 800247}
\textbf{Source problem code:} catalogue entry
\subsection{Definitions and mathematical statement}
Let \(n\geq 1\) be an integer. Interpret a complete-circular domain in the standard balanced sense: \(D\subset\C^n\) is a nonempty connected open set such that \(\lambda z\in D\) whenever \(z\in D\) and \(|\lambda|\leq 1\). Assume that \(D\) is bounded and pseudoconvex. Its Minkowski function is
\[
h(z)=\inf\{t>0:z/t\in D\},\qquad z\in\C^n,
\]
so that \(D=\{z\in\C^n:h(z)<1\}\). No continuity of \(h\) is assumed in the main question.

Let \(\Delta=\{\zeta\in\C:|\zeta|<1\}\), with Poincar\'e distance
\[
\rho(\zeta,\eta)=\operatorname{artanh}\left|\frac{\zeta-\eta}{1-\overline\eta\zeta}\right|.
\]
For definiteness, use the Carath\'eodory distance
\[
c_D(z,w)=\sup_{f\in\mathcal O(D,\Delta)}\rho(f(z),f(w)).
\]
Let \(k_D\) denote the Kobayashi distance, namely the largest pseudodistance on \(D\) for which every holomorphic map from \(\Delta\) to \(D\) decreases distance from \(\rho\).

To define the Bergman distance, put
\[
A^2(D)=\left\{f\in\mathcal O(D):\int_D|f|^2\,\mathrm dV<\infty\right\},\qquad
K_D(z,z)=\sup_{\substack{f\in A^2(D)\\ \|f\|_{L^2(D)}\leq 1}}|f(z)|^2.
\]
The Bergman metric and its induced path distance are
\[
\beta_D(z;v)^2=\sum_{j,\ell=1}^n
\frac{\partial^2\log K_D(z,z)}{\partial z_j\,\partial\overline z_\ell}
 v_j\overline v_\ell,
\qquad
b_D(z,w)=\inf_\gamma\int_0^1\beta_D(\gamma(t);\gamma'(t))\,\mathrm dt,
\]
where the infimum is over piecewise continuously differentiable curves in \(D\) joining \(z\) to \(w\). For \(d_D\in\{c_D,b_D,k_D\}\), completeness means that every \(d_D\)-Cauchy sequence converges in \(d_D\) to a point of \(D\).

\textbf{Problem (P. Pflug; source Problem 46).} Characterize, by properties of the Minkowski function \(h\), when \(D\) is Carath\'eodory complete, when it is Bergman complete, and when it is Kobayashi complete. A characterization is sought separately for each of these three notions, with necessary and sufficient conditions valid for the stated class of domains.

\textbf{Supplied historical qualifications.} The source reports the following as known: (a) continuity of \(h\) implies Bergman completeness of \(D\); (b) for \(n\geq 3\), continuity of \(h\) does not imply Kobayashi completeness of \(D\). These are assertions supplied with the 1995 problem, not independently verified results in this draft. The source does not specify whether Carath\'eodory completeness uses \(c_D\) as defined above or the path distance induced by the infinitesimal Carath\'eodory metric; that convention remains to be checked.
\subsection{Short English statement}
For a bounded pseudoconvex balanced domain, find necessary and sufficient conditions on its Minkowski function for completeness in each of the Carath\'eodory, Bergman, and Kobayashi distances. The source reports that continuity suffices for Bergman completeness but does not suffice for Kobayashi completeness in dimensions \(n\geq 3\). The precise Carath\'eodory convention requires review.
\subsection{Sources}
\begin{itemize}
\item Local UnsolvedMath registry, supplied complete source record: permanent ID \(800247\), exact problem code \texttt{catalogue entry}, title \emph{Several Complex Variables Problem List --- Problem 46}; upstream revision \texttt{d371ee917ac9647fbb06cd6fc529deec6f3dda82}. The registry labels the problem \texttt{open}; this is a supplied status, not an independently verified current assessment.
\item Buzzard, Fornaess, and Krantz, \emph{The Several Complex Variables Problem List} (1995), Problem 46, \url{https://arxiv.org/abs/math/9509202}. The supplied record attributes the individual problem to P. Pflug and gives the contact \texttt{pflugvec@dosuni1.bitnet}. The registry's \texttt{proposed\_by} field names Buzzard, Fornaess, and Krantz, while the extracted problem itself is signed by Pflug.
\item Supplied provenance: accessed 2026-10-16; extraction method \texttt{arxiv\_latex}. This access date and the historical assertions are supplied provenance, not evidence of independent source consultation or verification by this draft. No separate references for the two reported known results are supplied.
\end{itemize}
\end{document}
